\documentclass[aspectratio=169]{beamer}

\usetheme{Madrid}
\usecolortheme{seahorse}

\usepackage[utf8]{inputenc}
\usepackage{amsmath, amsfonts, amssymb}
\usepackage{booktabs}
\usepackage{array}
\usepackage{tikz}

\definecolor{AccentBlue}{RGB}{0,92,184}
\setbeamercolor{structure}{fg=AccentBlue}
\setbeamercolor{title}{fg=white,bg=AccentBlue}
\setbeamercolor{frametitle}{fg=AccentBlue}

\setbeamertemplate{navigation symbols}{}

% Placeholder macro for future figures/diagrams
\newcommand{\figplaceholder}[1]{%
  \begin{center}
    \begin{tikzpicture}
      \draw[dashed] (0,0) rectangle (8,4);
      \node at (4,2) {#1};
    \end{tikzpicture}
  \end{center}
}



\title[Contemporary Economic Issues]{What Are Contemporary Economic Issues?}
\subtitle{ECON 204 -- Contemporary Economic Issues}
\author{Muhebullah Karimzada}
%\institute{Winter 2026}
\date{Winter 2026}

\begin{document}

%------------------------------------------------
\begin{frame}
  %\transfade
  \titlepage
\end{frame}


%================================================
\section{Overview}

%------------------------------------------------
\begin{frame}{What Are Contemporary Economic Issues?}

  Contemporary economic issues are the main challenges that shape how people live and how
  economies function today. They are not abstract: they show up in news reports, social
  media discussions, political debates, and personal experiences.

  \vspace{0.4cm}

  \begin{itemize}
    \item When we talk about contemporary economic issues, we mean the main challenges that shape
    how people live and how economies function today.
    \item These topics are closely linked to everyday experiences of households, firms, and governments.
    \item Economics gives us a way to think about these issues with structure and evidence, rather than
    relying only on slogans or opinions.
  \end{itemize}
\end{frame}

%------------------------------------------------
\begin{frame}{Examples of Contemporary Economic Issues}

  Contemporary economic issues include topics such as:
  \begin{itemize}
    \item Rising cost of living and inflation
    \item Housing affordability and urban pressures
    \item Changes in jobs and wages due to technology
    \item Global trade tensions and supply chain disruptions
    \item Climate change and environmental policy
  \end{itemize}

  \vspace{0.3cm}
  These examples will be revisited throughout the course in more detail.
\end{frame}

%================================================
\section{Economics as a Way of Thinking}

%------------------------------------------------
\begin{frame}{Economics as a Way of Thinking}
  \transfade
  Economics is often defined as the study of how societies use limited resources to satisfy
  unlimited wants.

  \vspace{0.3cm}
  A useful way to think about economics is as a \textbf{way of thinking}, not just a set of formulas.

  \begin{itemize}
    \item Focus on how individuals, firms, and governments make decisions under constraints.
    \item Emphasis on trade offs, incentives, and interactions in markets and institutions.
    \item Use of models and data to organize our thinking and test ideas.
  \end{itemize}
\end{frame}

%------------------------------------------------
\subsection{Scarcity and Choice}

\begin{frame}{Scarcity and Choice}
  \transfade
  The starting point in economics is \textbf{scarcity}.

  \begin{itemize}
    \item Resources such as time, income, labour, land, and natural resources are limited.
    \item At the same time, people have many wants and needs.
    \item Because we cannot have everything at once, we must choose.
  \end{itemize}

  \vspace{0.3cm}
  Scarcity forces choice, and choice implies \textbf{trade offs}. This is a central idea in every
  contemporary economic issue.
\end{frame}

%------------------------------------------------
\begin{frame}{Examples of Scarcity and Choice}
  \transfade
  \begin{itemize}
    \item A student has 24 hours in a day. Time spent working a part-time job cannot be spent
    studying or relaxing.
    \item A government has a limited budget. Money allocated to healthcare cannot be used for
    schools or infrastructure at the same time.
    \item A piece of land in a city can be used for housing, a park, a shopping centre, or a school,
    but not all of these at once.
  \end{itemize}


\end{frame}

%------------------------------------------------
\subsection{Opportunity Cost}

\begin{frame}{Opportunity Cost}
  \transfade
  \textbf{Opportunity cost} is the value of the next best alternative that is given up when a decision
  is made. It is not simply about money; it is about what we could have done instead.

  \vspace{0.3cm}
  \begin{itemize}
    \item Every choice involves giving up something else.
    \item Comparing opportunity costs helps us rank alternatives.
    \item Policy debates often hinge on recognizing the opportunity cost of using resources in one
    way rather than another.
  \end{itemize}
\end{frame}

%------------------------------------------------
\begin{frame}{Examples of Opportunity Cost}
  \transfade
  \begin{itemize}
    \item If you spend an evening studying, your opportunity cost might be the movie night with
    friends that you miss.
    \item If a city chooses to build a new sports arena, the opportunity cost might be the public
    transit improvements that are delayed.
    \item If a government subsidizes one industry, the opportunity cost may be opportunities in
    other sectors that do not receive support.
  \end{itemize}

  \vspace{0.3cm}
  In modern policy debates, opportunity cost is everywhere. For example, when a government
  introduces a new housing subsidy, it needs to think not only about the direct cost of the
  program but also about what else could have been done with the same funds.
\end{frame}

%------------------------------------------------
\subsection{Incentives and Behaviour}

\begin{frame}{Incentives and Behaviour}
  \transfade
  Another key idea is that \textbf{incentives matter}. Incentives are rewards or penalties that
  encourage people to behave in certain ways.

  \begin{itemize}
    \item Prices, wages, taxes, and subsidies are all examples of economic incentives.
    \item Non-monetary incentives include social approval, career opportunities, or legal penalties.
    \item Hubbard et al.\ emphasize that when incentives change, behaviour tends to change as well.
  \end{itemize}
\end{frame}

%------------------------------------------------
\begin{frame}{Examples of Incentives}
  \transfade
  \begin{itemize}
    \item Higher wages in a sector attract more workers to that sector.
    \item A tax on cigarettes gives smokers a financial reason to smoke less.
    \item A subsidy for solar panels makes it more attractive for households to invest in renewable
    energy.
  \end{itemize}

  \vspace{0.3cm}
  Contemporary issues often involve complicated incentive structures:
  \begin{itemize}
    \item If unemployment benefits are too low, people may face severe hardship.
    \item If benefits are too generous, some may have less incentive to search actively for work.
    \item Good policy design aims to balance support and incentives.
  \end{itemize}
\end{frame}

%------------------------------------------------
\subsection{Positive and Normative Economics}

\begin{frame}{Positive and Normative Economics}
  \transfade
  Economics uses two kinds of statements:

  \vspace{0.2cm}
  \begin{itemize}
    \item \textbf{Positive statements} describe the world as it is. They can be checked against data.
          \begin{itemize}
            \item Example: ``An increase in the minimum wage reduced employment for teenagers in
            this region.''
          \end{itemize}
    \item \textbf{Normative statements} express judgments about what ought to be.
          \begin{itemize}
            \item Example: ``The minimum wage should be higher to ensure a decent standard of living.''
          \end{itemize}
  \end{itemize}
\end{frame}

%------------------------------------------------
\begin{frame}{Positive vs. Normative in Policy Debates}
  \transfade
  \begin{itemize}
    \item Real policy debates always mix positive and normative elements.
    \item Understanding the data and mechanisms (positive analysis) does not tell us what society
    should choose, but it clarifies the trade offs.
    \item Normative choices reflect ethical, social, and political values.
    \item Hubbard et al.\ stress that economists aim to separate ``what is'' from ``what should be'', even
    though both are important in public discussions.
  \end{itemize}
\end{frame}

%================================================
\section{What Makes an Issue ``Contemporary''?}

%------------------------------------------------
\begin{frame}{What Makes an Issue ``Contemporary''?}
  \transfade
  Many economic topics are long standing. An issue becomes \textbf{contemporary} when it is especially
  relevant to the current period.

  \vspace{0.3cm}
  Common features of contemporary issues:
  \begin{itemize}
    \item Real-time impact
    \item Uncertainty and risk
    \item Breaks with past patterns
    \item Fairness and distribution
  \end{itemize}
\end{frame}

%------------------------------------------------
\subsection{Real-Time Impact}

\begin{frame}{Real-Time Impact}
  \transfade
  A contemporary issue affects people directly now, or in the near future.

  \begin{itemize}
    \item A sudden rise in grocery prices is felt immediately at the checkout counter.
    \item A jump in rent shows up in monthly budgets and can force people to move.
    \item A sharp increase in interest rates quickly changes mortgage payments and business
    borrowing costs.
  \end{itemize}

  \vspace{0.3cm}
  Because the impact is real and immediate, these issues tend to receive attention in media
  and politics. Students can often connect these topics with their own experience.
\end{frame}

%------------------------------------------------
\subsection{Uncertainty and Risk}

\begin{frame}{Uncertainty and Risk}
  \transfade
  Contemporary issues often involve significant \textbf{uncertainty}. People are unsure how things will develop.

  \begin{itemize}
    \item Will artificial intelligence create more jobs than it destroys, or the reverse?
    \item How severe will climate-related damage be in 20 or 50 years?
    \item Will a trade conflict between two large countries escalate or be resolved?
  \end{itemize}

  \vspace{0.3cm}
  Economists pay attention to \textbf{expectations}. Expectations about future inflation, taxes, or
  regulations can influence decisions today.
\end{frame}

%------------------------------------------------
\begin{frame}{Expectations and Decisions}
  \transfade
  \begin{itemize}
    \item If people expect prices to rise quickly, they may try to buy sooner or ask for higher wages now.
    \item If firms expect higher future taxes or regulation, they may reduce investment today.
    \item In Hubbard et al., expectations are crucial in understanding interest rates, consumption,
    and investment behaviour over time.
  \end{itemize}
\end{frame}

%------------------------------------------------
\subsection{Breaks with Past Patterns}

\begin{frame}{Breaks with Past Patterns}
  \transfade
  Some issues are contemporary because they mark a clear shift away from historical patterns.

  \begin{itemize}
    \item Retail activity moving from physical stores to online platforms.
    \item The rise of gig work and contract positions relative to traditional full-time jobs.
    \item Global supply chains replacing earlier patterns of local production.
  \end{itemize}

  \vspace{0.3cm}
  When old structures change, institutions such as labour laws, tax systems, and social
  programs may no longer fit as well as before, creating pressure for policy reform and new ideas.
\end{frame}

%------------------------------------------------
\subsection{Fairness and Distribution}

\begin{frame}{Fairness and Distribution}
  \transfade
  Many contemporary issues are linked to questions of fairness and distribution:

  \begin{itemize}
    \item Who gains and who loses when trade expands?
    \item Who benefits most from new technologies?
    \item How are the costs of climate change shared across regions and generations?
  \end{itemize}

  \vspace{0.3cm}
  Economists measure inequality, study social mobility, and examine how different groups
  are affected by policy changes. Views on what is ``fair'' can differ, but understanding the
  distributional effects is essential for informed debate.
\end{frame}

%================================================
\section{Examples of Contemporary Economic Issues}

%------------------------------------------------
\subsection{Inflation and the Cost of Living}

\begin{frame}{Inflation and the Cost of Living}
  \transfade
  \textbf{Inflation} is the rate at which the average level of prices in the economy increases over time.

  \vspace{0.2cm}
  If the price level is $P_t$ at time $t$ and $P_{t+1}$ at time $t+1$, a simple inflation rate can be written as
  \[
    \pi = \frac{P_{t+1} - P_t}{P_t} \times 100\%.
  \]

  \vspace{0.2cm}
  \begin{itemize}
    \item When inflation is high, the purchasing power of money falls.
    \item Households may feel their paycheques do not go as far as before.
    \item Uncertainty about prices can make planning more difficult for firms and families.
  \end{itemize}
\end{frame}

%------------------------------------------------
\begin{frame}{Recent Inflation Experiences}
  \transfade
  Recent experience in many countries has included higher inflation after periods of stability.

  \begin{itemize}
    \item Strong demand following re-openings after the pandemic.
    \item Supply chain disruptions and shipping bottlenecks.
    \item Higher energy prices and other cost shocks.
  \end{itemize}

  \vspace{0.3cm}
  Central banks respond with \textbf{monetary policy}, often by raising or lowering interest rates
  to influence spending and inflation. This creates trade offs between controlling inflation and
  supporting employment and growth, a key theme in Hubbard et al.
\end{frame}

%------------------------------------------------
\subsection{Housing Affordability}

\begin{frame}{Housing Affordability}
  \transfade
  Housing is a basic need and also a major component of household wealth. In many cities,
  housing prices and rents have risen much faster than incomes.

  \vspace{0.2cm}
  Key ideas:
  \begin{itemize}
    \item \textbf{Demand factors}: population growth, migration to cities, higher incomes, and investor
    demand all increase the demand for housing.
    \item \textbf{Supply factors}: zoning rules, land scarcity, construction costs, and approval delays
    can limit how quickly new housing is built.
    \item \textbf{Interest rates}: lower interest rates reduce mortgage payments and can push prices
    up, while higher rates have the opposite effect.
  \end{itemize}
\end{frame}

%------------------------------------------------
\begin{frame}{Housing Affordability: Distributional Effects}
  \transfade
  \begin{itemize}
    \item Existing homeowners may see significant gains in wealth.
    \item Young households and renters may feel locked out of the market.
    \item Households may have to move farther from work or accept crowded conditions.
  \end{itemize}

  \vspace{0.3cm}
  Policy tools include:
  \begin{itemize}
    \item Changes in zoning and incentives for construction.
    \item Social housing programs and property tax reforms.
    \item Regulations on short-term rentals.
  \end{itemize}
\end{frame}

\subsection{Labour Markets and Technology}

\begin{frame}{Labour Markets and Technology}
  \transfade
  Technology has always influenced work, but the pace of change in recent decades has been
  particularly rapid. Automation, robotics, digital platforms, and artificial intelligence are
  transforming many occupations.

  \vspace{0.2cm}
  Important patterns:
  \begin{itemize}
    \item Machines and software can replace some routine tasks, such as data entry or simple
    assembly line operations.
    \item New tasks emerge that require higher skill levels in data analysis, programming, and
    creative problem solving.
    \item Work arrangements shift toward remote work, freelancing, and gig platforms.
  \end{itemize}
\end{frame}

%------------------------------------------------
\begin{frame}{Questions About Technology and Work}
  \transfade
  These shifts raise key questions:

  \begin{itemize}
    \item Will technology create more jobs than it destroys?
    \item How can education and training systems help workers transition to new roles?
    \item How should labour laws adapt to platform work and non-standard contracts?
  \end{itemize}

  \vspace{0.3cm}
  Economists study these questions using data on wages, employment, and skill requirements,
  and build models to see how shocks to technology affect different groups of workers.
\end{frame}

%------------------------------------------------
\subsection{Global Trade and Supply Chains}

\begin{frame}{Global Trade and Supply Chains}
  \transfade
  Global trade no longer means simply ``Country A exports finished goods to Country B''.

  \vspace{0.2cm}
  Many products are made through complex supply chains involving multiple stages in several
  countries.

  \begin{itemize}
    \item A car may have parts made in several countries, assembled in another, and sold worldwide.
    \item A smartphone may involve design in one country, chip production in another, assembly
    elsewhere, and final retail sales in many markets.
  \end{itemize}
\end{frame}

%------------------------------------------------
\begin{frame}{Benefits and Risks of Global Supply Chains}
  \transfade
  \begin{itemize}
    \item Global interconnection brings gains from specialization and lower costs.
    \item However, it also creates vulnerabilities:
          \begin{itemize}
            \item A natural disaster or political conflict in one region can disrupt production globally.
            \item Trade tensions and tariffs can raise costs and reduce access to markets.
            \item Dependence on a small number of suppliers for crucial inputs (such as semiconductors)
            can be risky.
          \end{itemize}
    \item Recent events have led many governments and firms to reconsider how global their
    supply chains should be, and whether more production should be brought closer to home.
  \end{itemize}
\end{frame}

%------------------------------------------------
\subsection{Climate Change and Environmental Policy}

\begin{frame}{Climate Change and Externalities}
  \transfade
  Climate change is both an environmental and an economic issue. It affects agriculture,
  coastal areas, infrastructure, health, and many sectors of the economy.

  \vspace{0.2cm}
  Economists often use the concept of an \textbf{externality}. An externality exists when an action
  by one party imposes a cost or benefit on others that is not reflected in market prices.
\end{frame}

%------------------------------------------------
\begin{frame}{Climate Change as a Negative Externality}
  \transfade
  In the case of climate change:

  \begin{itemize}
    \item A firm burning fossil fuels pays for the fuel but does not pay directly for the damage
    caused by greenhouse gas emissions.
    \item Individuals driving cars or heating homes contribute emissions but do not see the full
    long-term costs in their immediate bills.
  \end{itemize}

  \vspace{0.3cm}
  Because the full social cost of emissions is not included in private decisions, markets alone
  will not produce the desired level of emissions reductions.
\end{frame}

%------------------------------------------------
\begin{frame}{Environmental Policy Tools}
  \transfade
  Policy tools include:
  \begin{itemize}
    \item Carbon taxes that put a price on emissions.
    \item Cap-and-trade systems that limit total emissions and create tradable permits.
    \item Regulations and standards for vehicles, buildings, and industrial processes.
    \item Subsidies for clean technologies such as renewable energy and energy-efficient equipment.
  \end{itemize}

  \vspace{0.3cm}
  These policies involve trade offs between environmental protection, energy costs, and
  industrial competitiveness.
\end{frame}

%================================================
\section{Why Contemporary Issues Matter}

%------------------------------------------------
\subsection{For Households}

\begin{frame}{Why Contemporary Issues Matter for Households}
  \transfade
  For households, contemporary economic issues affect:
  \begin{itemize}
    \item Job security and wages.
    \item The ability to afford housing, food, education, and healthcare.
    \item Savings decisions and retirement planning.
  \end{itemize}

  \vspace{0.3cm}
  Understanding these issues helps individuals make better financial choices, interpret news
  about the economy, and evaluate claims made by politicians and commentators.
\end{frame}

%------------------------------------------------
\subsection{For Firms}

\begin{frame}{Why Contemporary Issues Matter for Firms}
  \transfade
  For firms, contemporary issues influence:
  \begin{itemize}
    \item Costs of inputs and labour.
    \item Demand for products and services.
    \item Investment decisions and location choices.
    \item Exposure to risk from policy changes, trade conflicts, and technological disruption.
  \end{itemize}

  \vspace{0.3cm}
  Managers and entrepreneurs rely on economic reasoning and data to make decisions under
  uncertainty.
\end{frame}

%------------------------------------------------
\subsection{For Governments}

\begin{frame}{Why Contemporary Issues Matter for Governments}
  \transfade
  Governments must design policies that:
  \begin{itemize}
    \item Stabilize the macroeconomy during recessions or periods of high inflation.
    \item Provide essential public services and infrastructure.
    \item Address inequality and support vulnerable groups.
    \item Respond to long-term challenges such as climate change and population aging.
  \end{itemize}

  \vspace{0.3cm}
  Because resources are limited, every new program has an opportunity cost. Economic analysis
  helps clarify these trade offs and predict likely effects of different policy options.
\end{frame}

%================================================
\section{How Economists Study Contemporary Issues}

%------------------------------------------------
\subsection{The Role of Data}

\begin{frame}{The Role of Data}
  \transfade
  Modern economics is highly empirical. Economists rely on data from statistical agencies,
  central banks, international organizations, and private sources.

  \vspace{0.2cm}
  Common data types:
  \begin{itemize}
    \item \textbf{Time series}: how key variables like GDP, unemployment, and inflation change over time.
    \item \textbf{Cross-sectional}: comparisons between individuals, firms, or regions at a point in time.
    \item \textbf{Panel data}: combinations of time and cross-section, following the same units across time.
  \end{itemize}
\end{frame}

%------------------------------------------------
\begin{frame}{Using Data in Practice}
  \transfade
  \begin{itemize}
    \item Students will often see graphs of data in class.
    \item Learning to read these graphs and understand basic trends is a key skill, even for
    students outside economics.
    \item Hubbard et al.\ use data throughout to connect theoretical models to real-world evidence.
  \end{itemize}
\end{frame}

%------------------------------------------------
\subsection{The Role of Models}

\begin{frame}{The Role of Models}
  \transfade
  Economic models provide simplified representations of reality. They focus on key relationships
  while ignoring some details.

  \vspace{0.2cm}
  Examples include:
  \begin{itemize}
    \item Supply and demand models for markets.
    \item Labour market models linking wages and employment to demand and supply for labour.
    \item Aggregate demand and aggregate supply models for the whole economy.
    \item Externality and public goods models for environmental and social issues.
  \end{itemize}
\end{frame}

%------------------------------------------------
\begin{frame}{Why Use Models?}
  \transfade
  \begin{itemize}
    \item Models are not perfect descriptions of the world, but they help organize thinking and
    clarify cause and effect.
    \item They allow economists to make predictions about how changes in policy or conditions
    are likely to affect outcomes.
    \item They provide a framework for interpreting data and distinguishing correlation from causation.
  \end{itemize}
\end{frame}

%------------------------------------------------
\subsection{Policy Evaluation}

\begin{frame}{Policy Evaluation}
  \transfade
  Economists use data and models to evaluate policy options. They ask questions such as:
  \begin{itemize}
    \item Does this policy increase total economic output?
    \item How are the gains and losses distributed across different groups?
    \item What are the short-run and long-run effects?
    \item Are there unintended consequences?
  \end{itemize}
\end{frame}

%------------------------------------------------
\begin{frame}{Economics and Democratic Choice}
  \transfade
  \begin{itemize}
    \item In practice, there is rarely a single ``correct'' policy.
    \item Different societies may make different choices, depending on their values and priorities.
    \item Economic analysis does not replace democratic decision making, but it provides important
    input by clarifying trade offs and likely outcomes.
  \end{itemize}
\end{frame}

%================================================
\section{Conclusion}

%------------------------------------------------
\begin{frame}{Conclusion}
  \transfade
  Contemporary economic issues touch nearly every aspect of daily life. They affect how we
  work, where we live, what we can afford, and how societies respond to big challenges such
  as climate change and technological disruption.

  \vspace{0.3cm}
  This lecture has highlighted three main ideas:
  \begin{itemize}
    \item Economics offers a structured way of thinking about scarce resources, trade offs, and
    incentives.
    \item An issue is contemporary when it is current, uncertain, connected to changing patterns,
    and often raises questions of fairness.
    \item Examples like inflation, housing, labour market change, trade, and climate policy show
    how economic reasoning can be applied to real-world problems.
  \end{itemize}
\end{frame}

%------------------------------------------------
\begin{frame}{Looking Ahead}
  \transfade
  \begin{itemize}
    \item In the rest of the course, each of these topics will be explored in more detail.
    \item We will develop more formal models and examine more specific evidence.
    \item The introductory ideas in this presentation will serve as a foundation for deeper work
    on contemporary economic issues.
  \end{itemize}
\end{frame}

%------------------------------------------------
\end{document}
